\section{DALL-E: Text-to-Image y Zero-Shot Transformation}

DALL-E coloca los text token y los image token en un mismo autoregressive model y aprende $p(\text{image}\mid\text{text})$. En la clase se usan interpolation, prompts complejos e image-to-image examples para mostrar su capacidad de composición. Los ejemplos son muy convincentes, pero también son los más fáciles de someter a cherry-pick; por eso la sample gallery debe tomarse como evidencia cualitativa y no como una estadística de la distribución completa.

\subsection{Interpolación y composición de conceptos}

Esta sección empieza con la interpolación de conceptos para observar cómo el condicionamiento textual desplaza la distribución visual. Al leer la figura hay que revisar al mismo tiempo la continuidad, la identidad de los objetos y las muestras fallidas, y no solo si la imagen resulta interesante. En la sección anterior, iGPT solo recibía un prefijo visual; aquí el lenguaje especifica los atributos objetivo, por lo que el cumplimiento de la condición y el realismo visual deben evaluarse por separado.

\lecturefigure{slide-25-dalle-interpolation.jpg}{DALL-E genera, mediante el condicionamiento textual, muestras de imágenes que varían de forma continua entre conceptos.}{Video oficial 29:10; Ramesh et al. 2021.}

\begin{knowledgebox}{Lectura de la figura: la interpolación muestra control por condición, no demuestra una geometría de conceptos perfecta}
Las prompt variation, como pasar de gato a perro, hacen que la apariencia de la salida cambie de forma continua, lo que indica que los token de texto pueden controlar la distribución visual. Hay que revisar las muestras fallidas, la prompt paraphrase, el object count y la spatial relation, para no fijarse solo en la grid más exitosa. La composición de conceptos también puede provenir de coocurrencias en los datos de entrenamiento y no de reglas abstractas.
\end{knowledgebox}

\subsection{Muestras diversas de texto a imagen}

A continuación se comparan varios grupos de prompts y samples para observar si diversity e instruction adherence se cumplen a la vez, y se registra el sesgo que puede introducir la gallery selection. Una misma fila responde «cuántos candidatos razonables puede generar el modelo»; la comparación entre filas responde «si al cambiar la condición textual cambian de manera predecible el objeto, el material y el estilo».

\lecturefigure{slide-26-dalle-examples.jpg}{Varios grupos de samples de DALL-E muestran la composición de condiciones de objeto, material, estilo y texto.}{Video oficial 30:13.}

\begin{knowledgebox}{Lectura de la figura: dentro de una fila se observa la diversity; entre filas, la instruction adherence}
Las muestras diversas de un mismo prompt muestran la distribution breadth; la comparación entre prompts distintos muestra si se respeta la condición. También hay que registrar de cuántos candidatos se seleccionó cada grid, si se usó un reranker y cuál es la tasa de fallos. Sin un selection protocol, la gallery no permite estimar la experiencia real del usuario.
\end{knowledgebox}

\subsection{Zero-shot image-to-image}

Después, la generación condicionada se extiende a image-to-image. El modelo puede volver a describir la imagen de entrada o transformarla hacia un concepto nuevo; aquí zero-shot significa que no se entrenó un mapping específico para esa transformation concreta, aunque los datos de preentrenamiento pueden contener relaciones similares.

\lecturefigure{slide-27-dalle-zero-shot-image2image.jpg}{La zero-shot transformation de DALL-E combina la entrada visual con la condición textual.}{Video oficial 30:33.}

\begin{knowledgebox}{Lectura de la figura: observar la conservación del contenido y el cambio de atributos}
Una transformation ideal conserva la identidad y la disposición, y al mismo tiempo cambia los atributos que indica el prompt. Ambos objetivos suelen entrar en conflicto: cambiar demasiado poco es no obedecer y cambiar demasiado es identity drift. La evaluación requiere tres tipos de metric, content preservation, condition adherence y perceptual quality, y no solo la estética.
\end{knowledgebox}

\lecturefigure{slide-28-dalle-more-image2image.jpg}{Más transformation examples exponen el alcance y la inestabilidad de la generalización composicional.}{Video oficial 31:33.}

\begin{knowledgebox}{Lectura de la figura: el valor de varios ejemplos está en encontrar los límites}
Al observar varios source/prompt, se ve qué objetos, posturas y fondos se conservan con facilidad y en cuáles ocurre colapso de modos o fuga de atributos. Unas cuantas muestras exitosas respaldan que «la capacidad existe», pero no permiten estimar que «la capacidad es confiable». En ingeniería conviene muestrear al azar, evaluar a ciegas y reportar la distribución de rechazos y fallos.
\end{knowledgebox}

\begin{warningbox}{«Parecerse» y «ser correcto» en la generación visual}
Una imagen no tiene una única respuesta correcta píxel por píxel, pero sigue sujeta a restricciones de hechos, identidad, derechos de autor, sesgo y seguridad. La perceptual quality no sustituye la evaluación de provenance, consent y policy; el espacio de salida abierto de los modelos generativos hace que la evaluación sea más difícil, no innecesaria.
\end{warningbox}

\subsection{Resumen de la sección}

DALL-E extiende el preentrenamiento autoregressive a la multimodal conditional generation. La gallery muestra que la capacidad existe; la confiabilidad, en cambio, requiere muestreo aleatorio, pruebas de composición y un selection protocol explícito.
